\title{DeepSeekMath: Pushing the Limits of Mathematical Reasoning in Open Language Models}

\begin{abstract}
Mathematical reasoning remains a substantial obstacle for language models because of its highly structured and intricate demands. In this work, we present DeepSeekMath~7B, which extends the pre-training of DeepSeek-Coder-Base-v1.5 7B by utilizing 120B tokens related to mathematics, compiled from Common Crawl, as well as supplementary natural language and code corpora. DeepSeekMath~7B attains a notable 51.7\% accuracy on the MATH benchmark at the competition level, doing so without resorting to external toolkits or ensemble approaches. Its performance is on par with models such as Gemini-Ultra and GPT-4. Using self-consistency sampling with 64 generations, DeepSeekMath~7B reaches a 60.9\% score on MATH. Two principal innovations account for DeepSeekMath~'s advanced mathematical reasoning: first, we leverage the vast pool of publicly accessible web content via a carefully constructed data filtering and selection framework; second, we propose Group Relative Policy Optimization (GRPO), a new variant of Proximal Policy Optimization (PPO), which not only improves reasoning capabilities in mathematics but also reduces the PPO memory footprint.
\end{abstract}

\section{Introduction}

Large language models (LLMs) have transformed mathematical reasoning within artificial intelligence, leading to marked improvements across quantitative reasoning benchmarks \citep{MATH} as well as geometry-focused evaluations \citep{trinh2024solving}. Furthermore, these systems have demonstrated value in supporting humans tackling intricate mathematical challenges \citep{tao}. Nevertheless, state-of-the-art models like GPT-4 \citep{gpt4} and Gemini-Ultra \citep{gemini} remain unavailable for public use, while existing open-source alternatives continue to lag significantly in terms of capability.

In this work, we present DeepSeekMath, a specialized language model tailored for mathematical domains, which demonstrates superior mathematical proficiency compared to existing open-source models and attains results comparable to GPT-4 on various academic evaluation tasks. Central to our approach is the development of the DeepSeekMath Corpus—a vast and meticulously curated pre-training dataset comprising 120B math-specific tokens. This resource is constructed from Common Crawl (CC) by deploying a fastText-based classifier \citep{joulin2016fasttext}. The classifier is initially trained with examples drawn from OpenWebMath \citep{openwebmath} as positives and a broad assortment of unrelated web pages as negatives. We then utilize the trained classifier to extract further positive samples from the CC; these candidates undergo rigorous refinement through manual annotation. The improved, annotated dataset is subsequently used to retrain the classifier for greater accuracy. Evaluation demonstrates the efficacy of our large-scale dataset, as evidenced by our model DeepSeekMath-Base 7B obtaining a score of 64.2\% on GSM8K \citep{gsm8k} and 36.2\% on the MATH competition dataset \citep{MATH}, surpassing Minerva 540B \citep{minerva}. Furthermore, given that DeepSeekMath Corpus supports multiple languages, we also observe marked gains on Chinese-language mathematical benchmarks \citep{wei2023cmath,agieval}. Our experience and results highlight promising directions in mathematical data collection for the research community, with ample potential for future progress.

DeepSeekMath-Base is built upon DeepSeek-Coder-Base-v1.5 7B \citep{deepseek-coder}, since initializing with a model pre-trained on code proves more effective than using a general-purpose LLM. Additionally, our findings suggest that training with mathematical data not only advances mathematical proficiency but also leads to improvements on benchmarks such as MMLU \citep{mmlu} and BBH \citep{bbh}. This demonstrates that the model's enhanced math skills come hand in hand with superior overall reasoning performance.

Following the pre-training phase, we fine-tune DeepSeekMath-Base for mathematical instructions using datasets that incorporate chain-of-thought \citep{cot}, program-of-thought \citep{pot,pal}, and tool-integrated reasoning techniques \citep{tora}. The resulting model, DeepSeekMath-Instruct 7B, surpasses other 7B models and achieves performance on par with open-source instruction-tuned models of 70B scale.

In addition, we present Group Relative Policy Optimization (GRPO), a novel reinforcement learning (RL) variant inspired by Proximal Policy Optimization (PPO) \citep{schulman2017proximal}. GRPO eliminates the need for a critic network, opting instead to derive baselines from aggregated group scores, which substantially decreases the computational demands during training. Leveraging only a selection of English-language instruction tuning data, GRPO demonstrates marked gains over the already competitive DeepSeekMath-Instruct, yielding improvements across both in-domain (GSM8K: 82.9\% $\rightarrow$ 88.2\%, MATH: 46.8\% $\rightarrow$ 51.7\%) and out-of-domain settings (e.g., CMATH: 84.6\% $\rightarrow$ 88.8\%) during RL fine-tuning. Furthermore, we introduce a unified framework that encompasses various approaches, such as Rejection Sampling Fine-Tuning (RFT) \citep{yuan2023scaling}, Direct Preference Optimization (DPO) \citep{dpo}, PPO, and GRPO, revealing that these methods can be viewed as either straightforward or simplified forms of RL. To further elucidate the nature of this paradigm, we carry out comprehensive empirical analyses—including comparisons between online and offline training, outcome versus process supervision, and single-turn versus multi-step RL. Finally, we articulate the reasons behind the enhanced performance of instruction-tuned models achieved by our RL strategies and propose avenues for further improving RL efficacy in this unified context.

\subsection{Contributions}
This work presents scalable mathematical pre-training methods and investigates reinforcement learning through comprehensive exploration and analysis.

\textbf{Math Pre-Training at Scale}

\begin{itemize}[topsep=0pt]
    \item This study demonstrates that the Common Crawl dataset, available to the public, harbors a wealth of information suitable for mathematical applications. Leveraging a carefully curated data selection workflow, we have developed the DeepSeekMath Corpus—an extensive, high-caliber dataset comprising 120B tokens derived from web sources specifically filtered for mathematical relevance. This corpus significantly surpasses previous resources, being nearly seven times greater than the collection of mathematical web pages used in Minerva \citep{minerva} and roughly nine times larger than the newly introduced OpenWebMath dataset \citep{openwebmath}.

\item The DeepSeekMath-Base 7B model, prior to any fine-tuning, attains a level of performance on par with Minerva 540B \citep{minerva}. This suggests that mathematical reasoning ability is influenced by factors beyond just model size. Notably, compact models that are pre-trained using well-curated datasets can also demonstrate impressive results.

\item We present results from our experiments on math training. Pre-training models on code enhances their performance on mathematical tasks, irrespective of tool usage. This provides partial evidence toward resolving the ongoing question: \textit{does code training enhance reasoning capabilities?} Our findings suggest that it does, specifically with regard to mathematical reasoning.

\item While it is standard practice to utilize arXiv papers for training—particularly within mathematics-focused works—this approach does not yield measurable gains on any of the mathematical benchmarks considered in this study.
\end{itemize}

\textbf{Exploration and Analysis of Reinforcement Learning}

\begin{itemize}[topsep=0pt]
    \item We propose Group Relative Policy Optimization (GRPO), a novel reinforcement learning method characterized by its efficiency and effectiveness. Unlike traditional algorithms that rely on a critic model, GRPO determines the baseline through group score estimation, which greatly reduces the computational and resource demands when compared to Proximal Policy Optimization (PPO).
    \item Our findings show that integrating GRPO with our instruction-tuned model, DeepSeekMath-Instruct, yields substantial gains in model performance utilizing only instruction-tuning datasets. Additionally, we notice improvements in the model's generalization capabilities across domains during reinforcement learning.
    \item We establish a cohesive framework for interpreting a range of approaches, encompassing RFT, DPO, PPO, and GRPO. Further, we perform a comprehensive series of experiments—including comparisons between online and offline training, evaluations of outcome versus process supervision, and assessments of single-turn versus iterative reinforcement learning—in order to thoroughly analyze the foundational components of this framework.
    \item Building upon this unified paradigm, we delve into the underlying mechanisms contributing to the success of reinforcement learning, and we outline several promising avenues for advancing the reinforcement learning of large language models.
\end{itemize}

\subsection{Summary of Evaluations and Metrics}

\begin{itemize}[topsep=0pt]
    \item \textbf{English and Chinese Mathematical Reasoning}: Our models are thoroughly evaluated using a variety of English and Chinese benchmarks that span mathematical tasks from elementary through collegiate levels. The English datasets examined are GSM8K \citep{gsm8k}, MATH \citep{MATH}, SAT \citep{llemma}, OCW Courses \citep{minerva}, and MMLU-STEM \citep{mmlu}. For Chinese, we utilize MGSM-zh \citep{mgsm}, CMATH \citep{wei2023cmath}, Gaokao-MathCloze \citep{agieval}, and Gaokao-MathQA \citep{agieval}. The evaluation focuses on the models' proficiency in independently producing written solutions and in solving the problems programmatically via Python.

On standard English evaluation tasks, DeepSeekMath-Base performs on par with the proprietary Minerva 540B model \citep{minerva}, and consistently outperforms all available open-source base models, such as Mistral 7B \citep{mistral} and Llemma-34B \citep{llemma}, irrespective of their use of math-specific pre-training, often achieving notably higher results. Remarkably, DeepSeekMath-Base leads on Chinese evaluation sets as well, which is plausibly attributed to our divergence from earlier studies \citep{minerva,llemma} that rely exclusively on English mathematical pre-training corpora; instead, we integrate high-quality non-English data. Through the incorporation of mathematical instruction fine-tuning and reinforcement learning techniques, DeepSeekMath-Instruct and DeepSeekMath-RL exhibit robust performance, surpassing 50\% accuracy on the highly challenging, competition-style MATH dataset—a milestone for open-source systems.

\item \textbf{Formal Mathematics}: DeepSeekMath-Base is assessed on its ability to convert informal mathematical statements to formal proofs, following the benchmark introduced in \citep{dsp_proof}, with the miniF2F dataset \citep{minif2f} and Isabelle proof assistant \citep{isabelle}. The results indicate that DeepSeekMath-Base achieves robust performance in few-shot autoformalization scenarios.

\item \textbf{Natural Language Understanding, Reasoning, and Code}: To holistically examine the breadth of the models' proficiency in general language understanding, reasoning, and programming, we assess DeepSeekMath-Base on the Massive Multitask Language Understanding (MMLU) benchmark \citep{mmlu}—which consists of 57 diverse multiple-choice assessments—as well as on BIG-Bench Hard (BBH) \citep{bbh}, a suite of 23 demanding tasks focused primarily on complex, multi-step reasoning. Additionally, we employ HumanEval \citep{codex} and MBPP \citep{mbpp} to evaluate code generation abilities. Pre-training on mathematical data provides tangible improvements in both language comprehension and reasoning outcomes.

\section{Math Pre-Training}

\subsection{Data Collection and Decontamination} This section details the methodology employed to assemble the DeepSeekMath Corpus utilizing data extracted from Common Crawl. As illustrated in Figure \ref{fig:data_collect}, we introduce a stepwise pipeline designed to incrementally curate a substantial collection of mathematical content from Common Crawl, beginning with an initial seed corpus (such as a compact, high-caliber set of mathematics datasets). Notably, this methodology is transferable and can be adapted to additional fields, including programming.

Our process begins with OpenWebMath \citep{openwebmath}, a curated repository of mathematical texts sourced from the web, as our primary seed dataset. Using this material, we construct and train a fastText model \citep{joulin2016fasttext} to help retrieve web pages that share characteristics with OpenWebMath content. For the model’s supervised training phase, we draw 500,000 positive samples from OpenWebMath and 500,000 negative samples from Common Crawl. We utilize an open-source fastText implementation\footnote{\url{https://fasttext.cc}}, setting hyperparameters as follows: a vector size of 256, a learning rate of 0.1, a maximum n-gram size of 3 words, a minimum word frequency threshold of 3, and three training epochs. In order to reduce redundancy in Common Crawl, we apply deduplication methods based on URL matching and similarity, yielding a processed dataset of approximately 40 billion HTML documents. The trained fastText model is then employed to extract web pages with mathematical relevance from the deduplicated Common Crawl. To ensure high content quality, candidate pages are ranked by their fastText-assigned scores, and only those at the top of the list are retained. The amount of preserved data is calibrated by conducting pre-training evaluations on sample sizes of the top 40B, 80B, 120B, and 160B tokens. For the initial stage, we opt to maintain the top 40B tokens.

Following the initial round of data gathering, a significant number of mathematical web pages remain missing. This shortfall is primarily due to the limited diversity present in the positive sample set used for training the fastText model. To remedy this, we seek out further mathematical web resources to broaden and diversify the seed corpus, with the aim of refining the fastText model's effectiveness. Our process begins by segmenting the entire Common Crawl dataset into mutually exclusive domains, each defined by sharing a common base URL across their constituent web pages. For each domain, we determine the proportion of web pages captured during the first data collection phase. Domains where more than 10\% of their pages have been collected are designated as math-related domains (for example, \url{mathoverflow.net}). We then conduct manual annotation of URLs within these math-associated domains that correspond to mathematical content (such as \url{mathoverflow.net/questions}). Any previously uncollected web pages linked to these URLs are subsequently incorporated into our seed corpus. By expanding our positive sample set in this way, we are able to train a more robust fastText model, thereby enhancing our ability to recover mathematical data in later collection phases. After conducting four such iterations, our efforts result in the accumulation of 35.5M mathematical web pages, comprising a total of 120B tokens. Notably, during the fourth round, we observe that almost 98\% of the data had already been obtained in the preceding iteration, prompting us to discontinue further data collection.

In order to prevent overlap with evaluation benchmarks, we adopt the methodology from \cite{deepseek-coder}, excluding web content that features either questions or solutions derived from English mathematical benchmarks like GSM8K \citep{gsm8k} and MATH \citep{MATH}, as well as Chinese resources such as CMATH \citep{wei2023cmath} and AGIEval \citep{agieval}. Specifically, we eliminate any text segment from the training data that contains a 10-gram phrase perfectly matching any sequence found within these benchmarks. For benchmark passages with lengths under 10 grams but at least 3 grams, we remove webpages based on exact substring matches, thus mitigating data leakage into our training set.

\subsection{Validating the Quality of the DeepSeekMath~Corpus}

We run pre-training experiments to investigate how the DeepSeekMath~Corpus is compared with the recently released math-training corpora:

\begin{itemize}[topsep=0pt]
    \item \textbf{MathPile} \citep{mathpile}: a collection of data drawn from diverse sources (totaling 8.9B tokens), including textbooks, Wikipedia, ProofWiki, CommonCrawl, StackExchange, and arXiv—with arXiv representing the dominant portion (exceeding 85\%) of the entire set;
    \item \textbf{OpenWebMath} \citep{openwebmath}: mathematical content identified and extracted from the CommonCrawl dataset, resulting in a corpus containing 13.6B tokens;
    \item \textbf{Proof-Pile-2} \citep{llemma}: a compilation of mathematical resources that brings together OpenWebMath, AlgebraicStack (comprising 10.3B tokens of mathematical programming code), and arXiv documents (amounting to 28.0B tokens). In our work using Proof-Pile-2, we employ the arXiv:Web:Code distribution of 2:4:1, as prescribed by \cite{llemma}.
\end{itemize}

\subsubsection{Training Setting}
\label{sec:quality-policy}
Mathematical training is performed on a broadly pre-trained language model containing 1.3B parameters, built upon the same architecture as the DeepSeek LLMs \citep{deepseek-llm}, hereafter referred to as DeepSeek-LLM 1.3B. For each distinct mathematical dataset, an individual model is trained across 150B tokens. All training procedures utilize the highly efficient HAI-LLM \citep{haillm} framework. Adhering to the methodologies established for DeepSeek LLMs, we employ the AdamW optimizer \citep{adamW} configured with $\beta_1=0.9$, $\beta_2=0.95$, and $\mathrm{weight\_decay}=0.1$. A multi-step learning rate schedule is adopted: the learning rate increases to its maximum after a 2,000-step warmup, is reduced to 31.6\% of its peak at 80\% of the total training steps, and is further decayed to 10.0\% of the peak at 90\% of the process. The learning rate peaks at 5.3e-4, with a training batch size of 4 million tokens and a maximum context length of 4,000 tokens.

\subsubsection{Evaluation Results}

\textbf{The DeepSeekMath~Corpus is of high quality, covers multilingual mathematical content, and is the largest in size.}

\begin{itemize}[topsep=0pt]
    \item \textbf{High-quality}:
    We assess the downstream capabilities across eight mathematical benchmarks utilizing few-shot chain-of-thought prompting \cite{cot}. Table \ref{tab:corpora_comparison} demonstrates a pronounced superiority in the outcomes of the model that has been trained with the DeepSeekMath~Corpus. Furthermore, Figure \ref{fig:corpora_comparisons} indicates that at 50B tokens—corresponding to one full epoch over Proof-Pile-2—the DeepSeekMath Corpus-trained model outperforms its Proof-Pile-2 counterpart, suggesting an overall higher data quality within the DeepSeekMath Corpus.
    \item \textbf{Multilingual}:
    Comprising diverse linguistic sources, the DeepSeekMath~Corpus primarily includes English and Chinese as its most frequent languages. According to Table \ref{tab:corpora_comparison}, the use of DeepSeekMath~Corpus for training substantially boosts mathematical reasoning skills in both English and Chinese. This stands in contrast to most existing math corpora, which focus predominantly on English and offer marginal—if any—gains for Chinese, sometimes even impeding performance in that language.
    \item \textbf{Large-scale}:
    The scale of the DeepSeekMath~Corpus greatly surpasses that of traditional mathematical corpora. As presented in Figure \ref{fig:corpora_comparisons}, DeepSeek-LLM 1.3B, trained with this extensive corpus, exhibits a sharper and more sustained improvement trajectory during learning. In comparison, standard baseline corpora are considerably smaller, with repeated training epochs, resulting in early stagnation of the model’s performance.
\end{itemize}

\subsection{Training and Evaluating DeepSeekMath-Base 7B}

This section presents DeepSeekMath-Base 7B, a foundational model distinguished by its robust mathematical reasoning skills. The model's initialization utilizes DeepSeek-Coder-Base-v1.5 7B \citep{deepseek-coder}, followed by training on a dataset totaling 500B tokens. The data composition comprises 56\% from the DeepSeekMath Corpus, 4\% derived from AlgebraicStack, 10\% sourced from arXiv, 20\% consisting of code from Github, with the final 10\% made up of natural language text from Common Crawl spanning both English and Chinese. Training is generally aligned with the configuration described in Section \ref{sec:quality-policy}, with the exceptions being a peak learning rate of 4.2e-4 and a token batch size set at 10M.

This work presents an in-depth evaluation of DeepSeekMath-Base 7B's mathematical proficiency, concentrating on its performance in generating self-contained mathematical solutions independently, solving mathematical questions with the assistance of tools, and engaging in formal theorem verification. Additionally, we extend our analysis to offer a broader overview of the base model’s competencies, examining its aptitude in natural language understanding, logical reasoning, and programming abilities.

\paragraph{Mathematical Problem Solving with Step-by-Step Reasoning}

We assess the mathematical problem-solving capabilities of DeepSeekMath-Base using few-shot chain-of-thought prompting \citep{cot} on eight distinct benchmarks in both English and Chinese. The selected benchmarks include tasks focused on quantitative reasoning, such as GSM8K \citep{gsm8k}, MATH \citep{MATH}, and CMATH \citep{wei2023cmath}, as well as multiple-choice question sets like MMLU-STEM \citep{mmlu} and Gaokao-MathQA \citep{agieval}. Together, these benchmarks span a broad array of mathematical disciplines, ranging in difficulty from elementary through college level.

Table \ref{tab:base_math_cot} demonstrates that DeepSeekMath-Base 7B achieves the highest scores across all eight evaluated benchmarks when compared to other open-source base models, such as the commonly adopted Mistral 7B \citep{mistral} and the newer Llemma 34B \citep{llemma}, which incorporates math-focused training using Proof-Pile-2 \citep{llemma}. Of particular significance is the model's performance on the competition-level MATH dataset, where DeepSeekMath-Base exceeds the next-best open-source base models by more than 10 percentage points and even surpasses Minerva 540B \citep{minerva}—a closed-source model based on PaLM \citep{palm}, with 77 times greater parameter count and additional pretraining on mathematical corpora.

\paragraph{Mathematical Problem Solving with Tool Use} We assess the capability of models to perform mathematical reasoning with external tools on the GSM8K and MATH benchmarks by employing few-shot program-of-thought prompting \citep{pot,pal}. For each question, the models are asked to generate a Python script, leveraging modules like \textit{math} and \textit{sympy} to handle complex calculations. The answer is determined by running the generated code. As indicated in Table \ref{tab:base_math_pot_proof}, DeepSeekMath-Base 7B achieves better results than the previous leading model, Llemma 34B.

\paragraph{Formal Mathematics}

The automation of formal proofs contributes significantly to guaranteeing both the correctness and dependability of mathematical proofs, while also improving productivity—a topic that has garnered increasing interest recently. In our study, we assess the capabilities of DeepSeekMath-Base 7B for the informal-to-formal proving challenge introduced in \citep{dsp_proof}, which involves producing a formal proof from an informal proposition, its formalization, and an informal proof outline. Our evaluation employs miniF2F \citep{minif2f}, a formal mathematics benchmark reflecting Olympiad-level problems, where for each item, a formal proof in Isabelle is generated using few-shot prompts. Consistent with the methodology in \cite{dsp_proof}, our approach uses language models to create proof sketches, and then the automated prover Sledgehammer \citep{sledgehammer} supplements any missing steps. According to Table \ref{tab:base_math_pot_proof}, DeepSeekMath-Base 7B achieves notable results in the automatic formalization of proofs.

\paragraph{Natural Language Understanding, Reasoning, and Code}

We assess the capabilities of models in natural language understanding using the MMLU benchmark \citep{mmlu}, in reasoning using BBH \citep{bbh}, and in programming through the HumanEval \citep{codex} and MBPP \citep{mbpp} tasks. According to Table \ref{tab:understanding_reasoning_code}, DeepSeekMath-Base 7B demonstrates substantial gains in both MMLU and BBH compared to its predecessor, DeepSeek-Coder-Base-v1.5 \citep{deepseek-coder}, underscoring the beneficial effects that mathematical training confers on understanding and reasoning abilities. Furthermore, by incorporating code tokens during continual training, DeepSeekMath-Base 7B successfully preserves the coding benchmark performance of DeepSeek-Coder-Base-v1.5. In sum, DeepSeekMath-Base 7B achieves notably stronger results than the general-purpose model Mistral 7B \citep{mistral} across all three evaluated benchmarks in reasoning and code generation.

\section{Supervised Fine-Tuning}

\subsection{SFT Data Curation}

A diverse instruction-tuning dataset for mathematics is developed, encompassing both English and Chinese questions spanning a range of domains and difficulty levels. Each problem is matched with solutions articulated in chain-of-thought (CoT) \citep{cot}, program-of-thought (PoT) \citep{pot,pal}, and formats that incorporate external tool usage \citep{tora}. Altogether, the dataset contains 776K training samples.

\begin{itemize}[topsep=0pt]
    \item \textbf{English mathematical datasets}: We provide tool-enhanced annotations for both GSM8K and MATH problem sets. Additionally, we utilize a selection of MathInstruct \citep{MathInstruct} data as well as the training portion of Lila-OOD \citep{lila}, where each question is accompanied by a CoT or PoT-based solution. The assembled English dataset encompasses a wide array of mathematical disciplines, including but not limited to algebra, probability, number theory, calculus, and geometry.
    \item \textbf{Chinese mathematical datasets}: Our dataset includes Chinese K-12 math problems that represent 76 distinct sub-domains, such as linear equations. Each problem is annotated with solutions employing both CoT and tool-integrated reasoning formats.
\end{itemize}

\subsection{Training and Evaluating DeepSeekMath-Instruct 7B}

This section presents DeepSeekMath-Instruct 7B, which is derived from DeepSeekMath-Base through mathematical instruction tuning. During the training process, examples are concatenated in random order until a context window of 4K tokens is filled. The model undergoes 500 training steps using a batch size of 256, with the learning rate fixed at 5e-5 throughout.

We evaluate models' mathematical performance both without and with tool use, on 4 quantitative reasoning benchmarks in English and Chinese. We benchmark our model against the leading models of the time:

\begin{itemize}[topsep=0pt]
    \item \textbf{Closed-source models} are comprised of the following: (1) the GPT series, particularly GPT-4 \citep{gpt4} and GPT-4 Code Interpreter~\footnote{\url{https://openai.com/blog/chatgpt-plugins\#\#code-interpreter}}, both recognized for their advanced capabilities, (2) Gemini Ultra and Pro \citep{gemini}, (3) Inflection-2 \citep{inflection-2}, (4) Grok-1~\footnote{\url{https://x.ai/model-card}}, as well as several newly launched models from Chinese firms, including (5) Baichuan-3~\footnote{\url{https://www.baichuan-ai.com}}, and (6) the most recent GLM-4~\footnote{\url{https://open.bigmodel.cn/dev/api\#glm-4}}

    from the GLM family \citep{glm}.

These general-purpose models have typically been subjected to various alignment processes.

\item \textbf{Open-source models} in this survey encompass a range of general architectures such as: (1) DeepSeek-LLM-Chat 67B \citep{deepseek-llm}, (2) Qwen 72B \citep{qwen}, (3) SeaLLM-v2 7B \citep{seallm}, and (4) ChatGLM3 6B \citep{chatglm3}. Additionally, the list features models exhibiting mathematical enhancements, including (5) InternLM2-Math 20B~\footnote{\url{https://github.com/InternLM/InternLM-Math}}, which is based on InternLM2 and further trained on mathematical tasks with subsequent instruction tuning; (6) Math-Shepherd-Mistral 7B, which leverages PPO training \citep{schulman2017proximal} on top of the Mistral 7B \citep{mistral} base using a process-supervised reward function; (7) the WizardMath collection \citep{wizardmath}, which advances the mathematical problem-solving abilities of both Mistral 7B and Llama-2 70B \citep{llama2} via the evolve-instruct paradigm (an instructional fine-tuning method relying on AI-generated tasks) and further PPO optimization, drawing mainly on datasets such as GSM8K and MATH; (8) MetaMath 70B \citep{metamath}, formed by fine-tuning Llama-2 70B on an augmented corpus comprised of GSM8K and MATH; (9) ToRA 34B \cite{tora}, which is derived from CodeLlama 34B and subsequently trained for tool-supported mathematical reasoning tasks; and (10) MAmmoTH 70B \citep{MathInstruct}, realized by instruction-tuning Llama-2 70B utilizing the MathInstruct collection.

Table \ref{tab:sft_rl_math} illustrates that, in evaluation scenarios without tool usage, DeepSeekMath-Instruct 7B achieves robust step-by-step reasoning capabilities. Specifically, on the MATH dataset, which targets competition-level problems, our model outperforms all existing open-source alternatives and most proprietary systems (such as Inflection-2 and Gemini Pro) by a margin of at least 9\% absolute. This performance advantage holds even when compared to models with significantly more parameters (for instance, Qwen 72B) or models optimized through math-specific reinforcement learning methods (such as WizardMath-v1.1 7B). Although DeepSeekMath-Instruct reaches the performance levels of proprietary Chinese models like GLM-4 and Baichuan-3 on the MATH benchmark, it does not yet exceed the scores of top-tier systems such as GPT-4 and Gemini Ultra.

In the evaluation scenario permitting the combination of natural language reasoning with programmatic tool utilization for problem solving, DeepSeekMath-Instruct 7B attains nearly 60\% accuracy on MATH, outperforming all previously released open-source models. Across the other benchmarks, our model achieves performance on par with DeepSeek-LLM-Chat 67B, the former leading model which possesses tenfold more parameters.

\subsection{Group Relative Policy Optimization} Reinforcement learning (RL) has demonstrated significant potential in enhancing the mathematical reasoning skills of LLMs following the Supervised Fine-Tuning (SFT) phase \citep{wang2023math,wizardmath}. Here, we present our novel and efficient RL method, Group Relative Policy Optimization (GRPO).

\subsubsection{From PPO to GRPO} Proximal Policy Optimization (PPO) \citep{schulman2017proximal} is an actor-critic RL algorithm that is widely used in the RL fine-tuning stage of LLMs \citep{ouyang2022training}. In particular, it optimizes LLMs by maximizing the following surrogate objective:

\begin{equation}
\footnotesize
    \mathcal{J}_{PPO}(\theta) = \mathbb{E}{[q \sim P(Q), o \sim \pi_{\theta_{old}}(O|q)]} \frac{1}{|o|} \sum_{t=1}^{|o|} \min \left[ \frac{\pi_\theta(o_{t} | q, o_{<t})}{\pi_{\theta_{old}}(o_{t} | q, o_{<t})} A_{t}, \text{clip} \left( \frac{\pi_\theta(o_{t} | q, o_{<t})}{\pi_{\theta_{old}}(o_{t} | q, o_{<t})}, 1 - \varepsilon, 1 + \varepsilon \right)  A_{t} \right] ,
\end{equation}

The above formulation for $\mathcal{J}_{PPO}(\theta)$ utilizes expectations over queries $q$ sampled from $P(Q)$ and sequences $o$ obtained from the prior policy $\pi_{\theta_{old}}(O|q)$. For each time step $t$ within the sequence, the objective is averaged by the sequence length $|o|$ and incorporates the minimum between the scaled advantage using the policy ratio and its clipped counterpart—capping the ratio between $1-\varepsilon$ and $1+\varepsilon$ to ensure updates remain within a constrained trust region.

Here, $\pi_{\theta}$ denotes the present policy model, while $\pi_{\theta_{old}}$ refers to the previous version. The variables $q$ and $o$ represent questions and corresponding outputs, respectively; these pairs are drawn from the question set and generated outputs of the earlier policy $\pi_{\theta_{old}}$. The hyper-parameter $\varepsilon$ is associated with the clipping operation in PPO, serving to stabilize the optimization process. The advantage term $A_t$ is calculated using Generalized Advantage Estimation (GAE) \citep{gae}, relying on future rewards $\{r_{\ge t}\}$ and an estimated value function $V_{\psi}$. As a result, PPO necessitates simultaneous training of a value function together with the policy. Additionally, to counteract excessive optimization of the reward model, it is customary to include a KL divergence penalty between the policy and a fixed reference model on each token within the reward signal, as described by \citep{ouyang2022training}, that is,

\begin{equation}
    r_{t} = r_\varphi(q, o_{\le t}) - \beta \log\frac{\pi_{\theta}(o_{t}|q, o_{<t})}{\pi_{ref}(o_{t}|q, o_{<t})},
\label{eq:PPO-reward}
\end{equation}

Here, $r_\varphi$ denotes the reward model, $\pi_{ref}$ refers to the reference model—typically set as the original SFT model—and $\beta$ represents the weight applied to the KL penalty.

As the value function employed in PPO is typically another model of comparable size as the policy model, it brings a substantial memory and computational burden. Additionally, during RL training, the value function is treated as a baseline in the calculation of the advantage for variance reduction. While in the LLM context, usually only the last token is assigned a reward score by the reward model, which may complicate the training of a value function that is accurate at each token. To address this, as shown in Figure \ref{fig:grpo}, we propose Group Relative Policy Optimization (GRPO), which obviates the need for additional value function approximation as in PPO, and instead uses the average reward of multiple sampled outputs, produced in response to the same question, as the baseline. More specifically, for each question $q$, GRPO samples a group of outputs $\{o_1, o_2, \cdots, o_G\}$  from the old policy  $\pi_{\theta_{old}}$  and then optimizes the policy model by maximizing the following objective:

\begin{equation}
\footnotesize
\begin{split}
    \mathcal{J}_{GRPO}(\theta) &= \mathbb{E}_{q \sim P(Q), \{o_i\}_{i=1}^G \sim \pi_{\theta_{old}}(O|q)} \Bigg[ \\
    & \frac{1}{G}\sum_{i=1}^G \frac{1}{|o_i|} \sum_{t=1}^{|o_i|} \bigg( \min \Big( \frac{\pi_\theta(o_{i,t} | q, o_{i,<t})}{\pi_{\theta_{old}}(o_{i,t} | q, o_{i,<t})} \hat{A}_{i,t},\, \text{clip}\Big( \frac{\pi_\theta(o_{i,t} | q, o_{i,<t})}{\pi_{\theta_{old}}(o_{i,t} | q, o_{i,<t})},\, 1 - \varepsilon,\, 1 + \varepsilon \Big) \hat{A}_{i,t} \Big) \\
    & \quad - \beta\, \mathbb{D}_{KL}\left[\pi_{\theta} || \pi_{ref}\right]\bigg)\Bigg],
\end{split}
\label{eq:GRPO-obj}
\end{equation}

where $\varepsilon$ and $\beta$ are hyper-parameters, and $\hat{A}_{i,t}$  is the advantage calculated based on relative rewards of the outputs inside each group only, which will be detailed in the following subsections. The group relative way that GRPO leverages to calculate the advantages, aligns well with the comparative nature of rewards models,  as reward models are typically trained on datasets of comparisons between outputs on the same question. Also note that, instead of adding KL penalty in the reward, GRPO regularizes by directly adding the KL divergence between the trained policy and the reference policy to the loss, avoiding complicating the calculation of  $\hat{A}_{i,t}$. And different from the KL penalty term used in (\ref{eq:PPO-reward}), we estimate the KL divergence with the following unbiased estimator \citep{kl_approx}:

\begin{equation}
\small
    \mathbb{D}_{KL}\left[\pi_{\theta} || \pi_{ref}\right] = \frac{\pi_{ref}(o_{i,t}|q,o_{i,<t})}{\pi_{\theta}(o_{i,t}|q,o_{i,<t})}- \log\frac{\pi_{ref}(o_{i,t}|q,o_{i,<t})}{\pi_{\theta}(o_{i,t}|q,o_{i,<t})} - 1,
\end{equation}

which is guaranteed to be positive.

\begin{algorithm}[t]
  \small
  \caption{Iterative Group Relative Policy Optimization}
  \textbf{Input} initial policy model $\pi_{\theta_{\text{init}}}$; reward models $r_\varphi$; task prompts $\mathcal{D}$; 
  hyperparameters $\varepsilon$, $\beta$, $\mu$
  \begin{algorithmic}[1]
    \State Set policy model $\pi_\theta \leftarrow \pi_{\theta_{\text{init}}}$
    \For{iteration = 1, \dots, I}
      \State Assign current policy $\pi_{ref} \leftarrow \pi_{\theta}$
      \For{step = 1, \dots, M}
        \State Draw a mini-batch $\mathcal{D}_b$ from $\mathcal{D}$
        \State Store previous policy $\pi_{\theta_{old}} \leftarrow \pi_{\theta}$
        \State For every question $q \in \mathcal{D}_b$, generate $G$ responses $\{o_i\}_{i=1}^G \sim \pi_{\theta_{old}} (\cdot \mid q)$
        \State For each $o_i$, evaluate rewards $\{r_i\}_{i=1}^G$ via $r_\varphi$
        \State Utilize group relative advantage estimation to derive $\hat{A}_{i,t}$ for the $t$-th token in $o_i$
        \For{GRPO iteration = 1, \dots, $\mu$}
          \State Perform a policy update on $\pi_\theta$ to increase the GRPO objective (Equation \ref{eq:GC-GRPO})
        \EndFor
      \EndFor 
      \State Enhance $r_\varphi$ using replay-based continual training
    \EndFor 
  \end{algorithmic}
  \textbf{Output} $\pi_\theta$
  \label{alg:iter-grpo}
\end{algorithm}

\subsubsection{Outcome Supervision RL with GRPO} Specifically, for a given question $q$, a collection of outputs $\{o_1, o_2, \cdots, o_G\}$ is drawn from the old policy $\pi_{\theta_{old}}$. Each output is evaluated using a reward model, producing a set of $G$ rewards $\mathbf{r} = \{r_1, r_2, \cdots, r_G\}$. These reward values are subsequently standardized by removing the group mean and scaling by the group standard deviation. In the outcome supervision framework, the normalized reward assigned at the conclusion of each output $o_i$ serves as the advantage for every token in that output, leading to $\hat{A}_{i, t} = \widetilde{r}_i = \frac{r_i- {\rm mean}(\mathbf{r})}{{\rm std}(\mathbf{r})}$. The policy parameters are then updated through maximization of the objective described in equation (\ref{eq:GRPO-obj}).

\subsubsection{Process Supervision RL with GRPO} 

Outcome supervision restricts feedback to a single reward given at the conclusion of each generated output, which can prove inadequate or inefficient for directing the policy during intricate mathematical reasoning. Building upon prior work \cite{wang2023math}, we also consider process supervision, in which feedback is delivered at the conclusion of every intermediate reasoning step. More precisely, consider a question $q$ and $G$ sampled completions $\{o_1, o_2, \cdots, o_G\}$. For each output sequence, a process reward model assigns scores at each step, resulting in a reward structure: $\mathbf{R} = \{ \{r_1^{index(1)},\cdots,r_1^{index(K_1)}\}, \cdots,  \{r_G^{index(1)},\cdots,r_G^{index(K_G)}\} \}$, where $index(j)$ denotes the token position at the end of the $j$-th step and $K_i$ indicates the total number of steps for the $i$-th output. All rewards are standardized using their global mean and standard deviation as follows: $\widetilde{r}_i^{index(j)} = \frac{r_i^{index(j)} - {\rm mean(\mathbf{R})}}{{\rm std(\mathbf{R})}}$.

For policy improvement, the process supervision regime computes the advantage associated with each token by summing the normalized rewards across all subsequent steps, that is, $\hat{A}_{i, t} = \sum_{index(j) \ge t} \widetilde{r}_i^{index(j)}$. The policy is then updated to maximize the process supervision objective as specified in equation (\ref{eq:GRPO-obj}).

\subsubsection{Iterative RL with GRPO} During reinforcement learning, as training advances, the previously trained reward model may become inadequate to effectively guide the updated policy model. Consequently, we investigate the approach of iterative RL using GRPO. As depicted in Algorithm \ref{alg:iter-grpo}, the iterative GRPO framework involves constructing updated datasets for the reward model, based on new samples obtained from the policy model. The reward model is further refined through a replay technique that integrates 10\% of prior samples. Subsequently, the policy model replaces the reference model and is updated continuously utilizing the improved reward model.

\subsection{Training and Evaluating DeepSeekMath-RL}

Our reinforcement learning (RL) experiments utilize DeepSeekMath-Instruct 7B as the base model. The RL training data comprises chain-of-thought-style questions from GSM8K and MATH, sourced from the SFT dataset, totaling approximately 144,000 questions. To isolate the influence of RL on tasks not represented during RL, we omit all other SFT questions. The construction of the reward model training dataset aligns with the methodology in \citep{wang2023math}. The initial reward model is trained using DeepSeekMath-Base 7B, employing a learning rate of $2\times10^{-5}$. For the GRPO algorithm, the policy model's learning rate is fixed at $1\times10^{-6}$, and we use a KL penalty coefficient of $0.04$. For each input question, $64$ candidate completions are generated. Maximum sequence length and batch size are both set to $1024$. Following each exploration round, the policy model undergoes only one update step. Evaluation of DeepSeekMath-RL 7B is carried out on the same benchmarks as DeepSeekMath-Instruct 7B. Here, chain-of-thought GSM8K and MATH problems are classified as in-domain, while all other benchmark tasks are treated as out-of-domain.

Table \ref{tab:sft_rl_math} presents the results for both open- and closed-source systems, evaluated using chain-of-thought and tool-augmented reasoning approaches on English and Chinese tasks. The findings are as follows: 1) DeepSeekMath-RL 7B achieves accuracy rates of 88.2\% on GSM8K and 51.7\% on MATH with chain-of-thought reasoning, outperforming all open-source competitors within the 7B to 70B parameter range and exceeding the results of most closed-source counterparts. 2) Notably, DeepSeekMath-RL 7B is trained exclusively with chain-of-thought formatted instruction tuning datasets from GSM8K and MATH, building upon the DeepSeekMath-Instruct 7B foundation. Even with this limited data variety, it consistently outperforms DeepSeekMath-Instruct 7B by every metric considered, illustrating the substantial gains delivered by reinforcement learning.

\section{Discussion}
Here, we present the results obtained from our pre-training procedures as well as our reinforcement learning (RL) experiments.

\subsection{Lessons Learnt in Pre-Training}

We begin by detailing our pre-training experience. Unless indicated otherwise, the default training configurations described in Section \ref{sec:quality-policy} are applied. It should be emphasized that references to the DeepSeekMath Corpus in this section pertain to an 89B-token dataset obtained from the second round of data collection.

\subsubsection{Code Training Benefits Mathematical Reasoning}

An often-cited but unproven hypothesis claims that exposure to code enhances reasoning skills. In this work, we provide preliminary evidence addressing this claim in the context of mathematics: code training enhances models’ mathematical reasoning capabilities, regardless of whether auxiliary tools are utilized.

To study how code training affects mathematical reasoning, we experimented with the following two-stage training and one-stage training settings:

\textbf{Two-Stage Training}

\begin{itemize}[topsep=0pt]
    \item \textbf{Code Training for 400B Tokens $\rightarrow$ Math Training for 150B Tokens}: The DeepSeek-LLM 1.3B model undergoes training on a dataset of 400B code tokens, which is subsequently followed by additional training on 150B tokens drawn from mathematical data;
    \item \textbf{General Training for 400B Tokens $\rightarrow$ Math Training for 150B Tokens}: For comparison, we conduct a parallel experiment where the first training stage employs 400B general tokens—sourced from a comprehensive general-domain corpus assembled by DeepSeek-AI—instead of code tokens, to assess whether code tokens offer a tangible benefit over general tokens in enhancing the model’s capability for mathematical reasoning.
\end{itemize}

\textbf{One-Stage Training}

\begin{itemize}[topsep=0pt]
    \item \textbf{Math Training Over 150B Tokens}: DeepSeek-LLM 1.3B undergoes training on a dataset comprising 150 billion mathematics-related tokens;
    \item \textbf{Joint Training Using 400B Code Tokens and 150B Math Tokens}: Sequencing math training after code training has been found to reduce coding capabilities. To address this, we explore the potential of integrating code and math tokens during a unified training phase to assess whether this approach simultaneously enhances mathematical reasoning and mitigates catastrophic forgetting.
\end{itemize}

\paragraph{Results}
The downstream performance across various training configurations is presented in Table \ref{tab:code-math} and Table \ref{tab:code-math-general-eval}.

Code training enhances program-aided mathematical reasoning capabilities in both two-stage and one-stage training scenarios. As illustrated in Table \ref{tab:code-math}, when utilizing the two-stage training approach, applying code training by itself notably improves performance in solving GSM8K and MATH problems through Python. Introducing math training in the subsequent stage brings about additional gains. Notably, within the one-stage training context, concurrently incorporating code tokens and math tokens addresses the problem of catastrophic forgetting, which can occur in two-stage training, while simultaneously supporting both coding tasks (see Table \ref{tab:code-math-general-eval}) and program-aided mathematical reasoning (see Table \ref{tab:code-math}).

Training on code data enhances mathematical reasoning abilities even in the absence of external tools. Within the two-stage training framework, the preliminary phase focused on code yields noticeable improvements on its own. Additionally, this stage facilitates more effective mathematical training afterwards, ultimately culminating in superior overall results. Conversely, when code tokens and math tokens are mixed during single-phase training, performance on mathematical reasoning tasks without tools diminishes. A possible explanation is that DeepSeek-LLM 1.3B, given its relatively small model size, is insufficiently capable of jointly internalizing both code and mathematical knowledge in a single stage.

\subsubsection{ArXiv Papers Seem Ineffective in Improving Mathematical Reasoning}

ArXiv papers are commonly included as a component of math pre-training data \citep{minerva,gpt-f,llemma,mathpile}. However, detailed analysis regarding their impact on mathematical reasoning has not been extensively conducted. Perhaps counter-intuitively, according to our experiments, arXiv papers seem ineffective in improving mathematical reasoning. We experiment with models of different sizes, including DeepSeek-LLM 1.3B and DeepSeek-Coder-Base-v1.5 7B \citep{deepseek-coder}, using arXiv corpora that underwent varied processing pipelines:

\begin{itemize}[topsep=0pt]
    \item \textbf{MathPile} \citep{mathpile}: a dataset containing 8.9 billion tokens, assembled using various heuristic methods for data cleaning and filtering, with scientific arXiv papers constituting more than 85\% of its content;
    \item \textbf{ArXiv-RedPajama} \citep{redpajama}: comprises the full collection of arXiv LaTeX documents, from which preambles, comments, macros, and bibliographies have been stripped, resulting in a total size of 28.0 billion tokens.
\end{itemize}

In these experiments, we independently train DeepSeek-LLM 1.3B on 150B tokens and DeepSeek-Coder-Base-v1.5 7B on 40B tokens derived from each arXiv dataset. The results indicate that exposure to arXiv articles alone does not contribute meaningfully to enhancing mathematical reasoning abilities. When the models are trained exclusively on arXiv data, there is an absence of significant progress—and in some cases, even a decline—across several mathematical evaluation tasks of varying difficulty considered in this work. The evaluation encompasses quantitative reasoning benchmarks such as GSM8K and MATH (Table \ref{tab:arxiv-cot}), multiple-choice exams including MMLU-STEM (Table \ref{tab:arxiv-cot}), as well as formal mathematics tests exemplified by miniF2F (Table \ref{tab:arxiv_atp}).

However, this conclusion has its limitations and should be taken with a grain of salt. We have not yet studied:

\begin{itemize}[topsep=0pt]
    \item The influence of arXiv tokens on other mathematics-oriented tasks beyond the present study, for example, transforming formal statements or proofs into informal expressions (the process of informalization);
    \item The outcomes of integrating arXiv tokens with different categories of data;
    \item The extent to which the advantages conferred by arXiv papers might persist or expand when utilized with larger-scale models.
\end{itemize}

Thus, further exploration is required, which we leave for future studies.

\subsection{Insights of Reinforcement Learning}
\subsubsection{Towards to a Unified Paradigm} Here, we introduce a comprehensive framework for examining various training approaches, including SFT, RFT, DPO, PPO, and GRPO, and present experimental investigations into the components of this unified framework. In a general form, the parameter gradient $\theta$ for a given training method can be expressed as:

\begin{equation}
    \nabla_{\theta}\mathcal{J}_{\textcolor{red}{\mathcal{A}}}(\theta) = \mathbb{E}[\underbrace{(q,o) \sim \textcolor{red}{\mathcal{D}}}_{Data \ Source}]\left( \frac{1}{|o|} \sum_{t=1}^{|o|}  \underbrace{GC_{{\mathcal{A}}}(q, o, t, \textcolor{red}{\pi_{{rf}}})}_{Gradient \ Coefficient}  \nabla_{\theta}\log \pi_{\theta}(o_t | q, o_{<t})\right).
\label{eq:objective}
\end{equation}

There exist three key components: 1) \textit{Data Source $\mathcal{D}$}, which determines the training data; 2) \textit{Reward Function $\pi_{{rf}}$}, which is the source of the training reward signal; 3) \textit{Algorithm $\mathcal{A}$}: which processes the training data and the reward signal to the gradient coefficient $GC$ that determines the magnitude of the penalty or reinforcement for the data. We analyze several representative methods based on such a unified paradigm:

\begin{itemize}
    \item \textbf{Supervised Fine-tuning (SFT)}: SFT adapts a pretrained model by further training it on datasets specifically curated and labeled by humans.
    \item \textbf{Rejection Sampling Fine-tuning (RFT)}: In RFT, the SFT-trained model is fine-tuned further using a subset of its generated responses, chosen according to their correctness in relation to SFT questions.
    \item \textbf{Direct Preference Optimization (DPO)}: DPO extends the SFT approach by employing a pairwise loss function during fine-tuning, leveraging additional samples derived from the SFT model’s outputs.
    \item \textbf{Online Rejection Sampling Fine-tuning (Online RFT)}: Unlike the offline approach in RFT, Online RFT continually updates the policy model—starting from the SFT initialization—by fine-tuning on augmented samples generated on-the-fly from the live policy model.
    \item \textbf{PPO/GRPO}: PPO/GRPO initializes its policy from the SFT model and enhances the policy iteratively with real-time samples generated by the current version of the policy model using reinforcement learning.
\end{itemize}

Table \ref{tab:data_policy} provides an overview of the elements comprising these methods. For an in-depth explanation of the derivation steps, consult Appendix \ref{app:analysis-rl}.

\paragraph{Observation about Data Source} The data sources can be classified into two types: online sampling and offline sampling. Online sampling refers to cases where training data are derived from the outputs generated by the current training policy model during exploration. Conversely, offline sampling refers to training data obtained from outputs sampled by the initial SFT model. Both RFT and DPO utilize the offline sampling method, whereas Online RFT and GRPO adopt an online sampling approach.

As illustrated in Figure \ref{fig:rl-analysis}, our results demonstrate that Online RFT achieves markedly better performance than RFT on both benchmark tasks. During the initial phase of training, Online RFT performs at a similar level to RFT; however, as training progresses, it establishes a clear lead, underscoring the effectiveness of online training methods. This pattern is to be expected: early in training, the actor and the SFT model remain quite similar, so the data generated through sampling shows only minimal divergence. As training advances, the actor's sampled data becomes increasingly distinct, making the benefits of continual, real-time data collection and training more pronounced.

\paragraph{Observation about Gradient Coefficient} The algorithm translates the reward signal into a gradient coefficient to facilitate updates to the model parameters. In our study, we categorize the reward function into 'Rule' and 'Model' components. The 'Rule' approach evaluates response quality strictly in terms of answer correctness, whereas the 'Model' involves training a reward model that assigns scores to responses. This reward model utilizes data labeled by the 'Rule' criterion. As shown in Equations \ref{eq:GC-RFT} and \ref{eq:GC-GRPO}, an important distinction emerges between GRPO and Online RFT: GRPO dynamically modulates its gradient coefficient in relation to the reward values assigned by the reward model, thus enabling selective reinforcement or penalization according to the graded feedback. Conversely, Online RFT does not incorporate this mechanism—it neither penalizes wrong responses nor distinguishes among correct responses, instead applying uniform positive reinforcement across all correct answers.

As depicted in Figure \ref{fig:rl-analysis}, GRPO outperforms online RFT, which underscores the effectiveness of modifying the coefficients for positive and negative gradients. Moreover, GRPO+PS achieves better results than GRPO+OS, suggesting that employing more granular, step-sensitive gradient coefficients yields advantages. Additionally, we examine the impact of iterative RL by performing two iterations in our experiments. Figure \ref{fig:iter-rl} reveals that incorporating iterative RL notably enhances performance, with the most pronounced gains observed during the initial iteration.

\subsubsection{Why RL Works?} In our study, we implement reinforcement learning utilizing a select portion of instruction tuning datasets and observe that it markedly improves model performance compared to instruction tuning alone. To probe the underlying reasons for the effectiveness of reinforcement learning, we assess both Pass@K and Maj@K accuracies for Instruct and RL models across two benchmark tasks. The results, visualized in Figure \ref{fig:maj-pass}, demonstrate that reinforcement learning boosts Maj@K accuracy, yet leaves Pass@K essentially unchanged. This suggests that reinforcement learning bolsters the overall output distribution by strengthening the probability of producing correct answers among the TopK predictions, rather than fundamentally advancing the model’s intrinsic abilities. In a similar vein, \citep{wang2023making} uncovered a \textbf{misalignment problem} in reasoning tasks for SFT models, noting that reasoning performance can be enhanced by implementing a variety of preference alignment methodologies \citep{yuan2023rrhf,song2023preference,wang2023making}.

\subsubsection{How to Achieve More Effective RL?}
\label{sec:effec_rl}
Our findings indicate that RL performs efficiently on mathematical reasoning problems. Additionally, we introduce a cohesive framework for interpreting various prominent training strategies. Under this framework, these strategies are interpreted as forms of either full-fledged or streamlined RL approaches. As detailed in Equation \ref{eq:objective}, this framework highlights three principal elements: Data Source, Algorithm, and Reward Function. Below, we outline several possible avenues for further research concerning these elements.

\paragraph{Data Source} The data source forms the foundational basis for all training approaches. Within reinforcement learning (RL), this term denotes the collection of unlabeled questions paired with responses produced by sampling from the policy model. Throughout this work, we restrict ourselves to using only the questions introduced during the instruction tuning phase and employ basic nucleus sampling to generate outputs. We hypothesize that this limitation may underlie why our RL framework leads to performance gains predominantly in Maj@K. For future studies, we intend to apply our RL pipeline to question prompts outside the original training distribution, together with the adoption of \textbf{advanced sampling (decoding) strategies}, such as tree-search-based methods \citep{yao2023tree}. Furthermore, \textbf{efficient inference techniques} \citep{xia-etal-2023-speculative,leviathan2023fast,kwon2023efficient,xia2024unlocking}, which influence how effectively the policy model explores, are also crucial factors that warrant attention.

\paragraph{Algorithms} Algorithms utilize the data and corresponding reward signals to compute gradient coefficients, which are then used to adjust the model's parameters. According to Equation \ref{eq:objective}, current approaches universally exhibit a high level of \textbf{TRUST} in the reward function's feedback, modulating the conditional probability of specific tokens based on these signals. Nevertheless, the assumption that reward signals consistently offer accurate guidance does not always hold, particularly for highly intricate tasks. As an illustration, the PRM800K datasets \citep{lightman2023let}, despite extensive curation and annotation by skilled annotators, still demonstrate an error rate with about 20\% incorrect annotations\footnote{\url{https://github.com/openai/prm800k/issues/12\#issuecomment-1728491852}}. In light of this, our investigation will focus on reinforcement learning algorithms that maintain robustness even when confronted with noisy reward signals. We posit that leveraging \textbf{WEAK-TO-STRONG} \citep{burns2023weak} alignment techniques will yield transformative advancements for learning algorithms.

\paragraph{Reward Function} The reward function serves as the principal training signal in reinforcement learning, with the neural reward model typically playing this role. There are three critical areas for advancing reward models: 1) \textbf{Strategies to improve the generalization capability of the reward model.} To truly advance the fundamental abilities of large language models (LLMs), the reward model must generalize robustly to novel or out-of-distribution prompts as well as sophisticated generated responses; without this, reinforcement learning might only help LLMs maintain their existing output distribution instead of fostering substantive improvements; 2) \textbf{Ways to capture and represent the reward model's uncertainty.} Accurately modeling uncertainty can serve as a valuable connection between the limitations of existing reward models and the potential of learning paradigms that transition from weak to strong models; 3) \textbf{Approaches for building high-quality process reward models efficiently,} which are capable of generating nuanced, fine-grained feedback during the reasoning process \citep{lightman2023let,wang2023math}.

\section{Conclusion, Limitation, and Future Work}  
In this work, we introduce DeepSeekMath, which achieves superior results compared to all other open-source models on the MATH benchmark at the competition level and rivals the effectiveness of proprietary alternatives. DeepSeekMath~is built upon DeepSeek-Coder-v1.5 7B, undergoing continuous training over 500B tokens, with 120B of these drawn from Common Crawl as mathematical content. Our comprehensive ablation analysis indicates that web sources contribute valuable and high-quality mathematical datasets, whereas the inclusion of arXiv data does not yield the anticipated benefits. We propose Group Relative Policy Optimization (GRPO)—an adaptation of Proximal Policy Optimization (PPO)—demonstrating its ability to strengthen mathematical reasoning with improved memory efficiency. Experimental results confirm that even after attaining strong benchmark results, the use of GRPO further enhances DeepSeekMath-Instruct 7B. Additionally, we present a unified perspective to conceptualize a range of methodologies and highlight several avenues for advancing reinforcement learning effectiveness.

While DeepSeekMath~demonstrates strong performance on quantitative reasoning tasks, its abilities in geometry and theorem-proof areas are noticeably limited compared to proprietary models. For example, in our preliminary experiments, the model struggled with questions pertaining to triangles and ellipses, which could suggest that there was a bias in the data used during the pre-training and fine-tuning stages. Furthermore, due to constraints in model size, DeepSeekMath~does not match GPT-4 in few-shot scenarios; unlike GPT-4, which benefits from few-shot context, DeepSeekMath~yields nearly identical results in both zero-shot and few-shot settings. Going forward, we plan to enhance our data engineering pipeline to assemble a higher-quality pre-training dataset. Additionally, we aim to investigate novel strategies (see Section \ref{sec:effec_rl}) for optimizing reinforcement learning in large language models.

\end{document}